\section{A signed first-return section with all endpoints}\label{sec:boundary}

Every periodic orbit not wholly in $B_0$ has a coordinate of absolute
value $R+1$ or $R+2$, by Proposition~\ref{prop:integer-box}. Hence it
meets the section
\begin{equation}\label{eq:section}
 A=\{(x,y)\in B: |y|\in\{R+1,R+2\}\}.
\end{equation}
The map $h_d$ commutes with $S$. For a pair of opposite section
points, use the unique representative with positive second coordinate:
\[
 [x,\lambda]=(x,R+\lambda),\qquad \lambda\in\{1,2\}.
\]
This notation is only a proof device. Signs on the return edges
retain the ordinary, nonquotiented dynamics.

A return edge
\[
 [x,\lambda]\ \xto{\tau,\eps}\ [x',\lambda']
\]
means that $\tau$ is the first positive return time to $A$ and
$h_d^\tau(x,R+\lambda)=\eps(x',R+\lambda')$, where
$\eps\in\{+1,-1\}$. If a trajectory leaves $B$ instead, it is
discarded from periodic classification.

\begin{lemma}[Ordinary lift of a signed return cycle]\label{lem:lift}
Let a primitive normalized section cycle have total time $T$ and
edge-sign product $\eps$. If $\eps=+1$, it lifts to two distinct
primitive $h_d$-cycles of length $T$. If $\eps=-1$, it lifts to one
primitive $h_d$-cycle of length $2T$.
\end{lemma}
\begin{proof}
Following the return edges once sends a representative $z$ to
$\eps z$. Since $A$ has no point fixed by $S$, a negative sign
requires two traversals to return to $z$, whereas a positive sign
requires one. An earlier ordinary return would imply an earlier
repeat of the normalized section point, contradicting its primitive
return cycle and the first-return property. In the positive case
the lifts through $z$ and $-z$ cannot coincide: such a coincidence
would give a negative return to the same normalized starting point
before the positive traversal closed. Thus the stated periods and
cycle counts are ordinary ones.
\end{proof}

We use the endpoint labels
\begin{equation}\label{eq:endpoints}
 L_c=[-R+c,1],\qquad U_c=[R+c,1],\qquad V_c=[R+c,2].
\end{equation}
Let $r=R\bmod6$, $\kappa=3\lfloor r/3\rfloor$, and
$e=(x-\kappa)\bmod6$. The phase adjustment in $e$ makes it possible
to combine pairs of radius classes in the next table.

\begin{lemma}[All generic inner returns]\label{lem:strips}
For $[x,1]$ with $-R+L\le x\le R+U$ and residue $e$ as above,
Table~\ref{tab:strips} gives its first return. Every edge has sign
$-1$ and normalized target $[x+\delta,\lambda']$.
\end{lemma}

\begin{center}
\begin{minipage}{\linewidth}
\centering
\begin{tabular}{ccccccc}
\toprule
$s$ & $e$ & $L$ & $U$ & $\tau$ & $\delta$ & $\lambda'$\\
\midrule
$0$ & $0,3$ & $0$ & $0$ & $2$ & $0$ & $1$\\
$0$ & $1$ & $2$ & $0$ & $10$ & $-2$ & $1$\\
$0$ & $2$ & $4$ & $0$ & $14$ & $-4$ & $1$\\
$0$ & $4,5$ & $0$ & $0$ & $2$ & $0$ & $2$\\
\midrule
$1$ & $0$ & $4$ & $2$ & $6$ & $-4$ & $1$\\
$1$ & $1$ & $3$ & $-1$ & $18$ & $0$ & $1$\\
$1$ & $2,5$ & $2$ & $2$ & $2$ & $-2$ & $1$\\
$1$ & $3,4$ & $2$ & $2$ & $2$ & $-2$ & $2$\\
\midrule
$2$ & $0,5$ & $1$ & $1$ & $2$ & $-1$ & $2$\\
$2$ & $1,4$ & $1$ & $1$ & $2$ & $-1$ & $1$\\
$2$ & $2$ & $1$ & $-3$ & $10$ & $3$ & $1$\\
$2$ & $3$ & $1$ & $-1$ & $6$ & $1$ & $1$\\
\bottomrule
\end{tabular}
\captionof{table}{Generic inner first returns; all signs are negative.}
\label{tab:strips}
\end{minipage}
\end{center}

\begin{proof}
Keep $R$ and $x$ as free affine variables, with their residues fixed.
Starting from $(x,R+1)$, the first step is $(R+1,a-x)$.
Use \eqref{eq:values} at the initial boundary coordinate and
$p\sine(y)$ at each subsequent core coordinate. Record the inequality
$-R\le y\le R$ at every intermediate second coordinate. The exact
intersection of its $x$ inequalities is the displayed interval.
All other coordinates encountered before the return are of the form
$\pm R+c$ or constants; their core inequalities hold on the relevant
radius progression already at its least radius.

More explicitly, for each of the six choices of $r$ and each
$e=0,\ldots,5$, substitute $x\equiv e+\kappa\pmod6$ into the
affine recurrence. Iterate until the second coordinate is identically
$\pm(R+\lambda')$. In all $36$ cases this occurs at the listed
$\tau$, the last pair is $-(x+\delta,R+\lambda')$, and the
intersected bounds have the stated offsets. At each earlier time the
recorded second coordinate is in $[-R,R]$, so this is a first return.
These steps are a finite identity check in two affine variables;
they verify the table for every $R$ in the progression simultaneously.
The least radii for $r=0,1,\ldots,5$ are respectively
$6,7,2,3,4,5$, which also cover the smallest allowed degrees.
\end{proof}

For an outer state $[x,2]$, the next second coordinate is
$2R+s-x$. Thus it leaves $B$ immediately if
\begin{equation}\label{eq:outer-threshold}
 x<R+s-2.
\end{equation}
Every other outer representative is one of the finitely many
$V_c$ with $s-2\le c\le2$. There is no unexamined long outer strip.

The inner endpoint list is likewise an exact complement, not a
bounded search near corners. For any row of Table~\ref{tab:strips},
the complement of its generic interval in $[-R-2,R+2]$ is
\begin{equation}\label{eq:complement}
 x=-R+c,\quad -2\le c\le L-1,
 \qquad\text{or}\qquad
 x=R+c,\quad U+1\le c\le2,
\end{equation}
filtered by that row's residue $e$. The generic real interval is
nonempty at the least allowed radius, since $2R+U-L\ge0$ there.
The two complementary intervals are consequently disjoint. Applying
\eqref{eq:complement} gives exactly six exceptional inner states in
each radius class. Together with the surviving outer states they
give the complete endpoint table below.

\begin{lemma}[Complete endpoint returns and escapes]\label{lem:endpoints}
All section representatives outside the generic inner strips either
escape immediately by \eqref{eq:outer-threshold} or occur as a start
in Table~\ref{tab:endpoints}. The table gives all nonescaping first
returns. The symbol $E$ denotes an escape path, whose precise short
transient is specified after the table.
\end{lemma}

{\small
\begin{longtable}{cllcc}
\caption{Complete exceptional inner and surviving outer rules.
The paired $L$ row for $s=1$ represents two states.}
\label{tab:endpoints}\\
\toprule
$s$ & Start & First normalized target & Time & Sign\\
\midrule
\endfirsthead
\multicolumn{5}{c}{Table \thetable\ continued}\\
\toprule
$s$ & Start & First normalized target & Time & Sign\\
\midrule
\endhead
\midrule
\multicolumn{5}{r}{Continued on next page}\\
\endfoot
\bottomrule
\endlastfoot
$0$ & $L_{-2}$ & $V_1$ & $1$ & $+$\\
$0$ & $L_{-1}$ & $U_1$ & $1$ & $+$\\
$0$ & $L_1$ & $U_{-1}$ & $7$ & $-$\\
$0$ & $L_2$ & $U_{-2}$ & $9$ & $+$\\
$0$ & $V_{-2}$ & $V_2$ & $1$ & $+$\\
$0$ & $V_{-1}$ & $U_2$ & $1$ & $+$\\
$0$ & $V_0$ & $E$ & --- & ---\\
$0$ & $U_1$ & $L_{-1}$ & $1$ & $-$\\
$0$ & $V_1$ & $L_1$ & $2$ & $-$\\
$0$ & $U_2$ & $E$ & --- & ---\\
$0$ & $V_2$ & $L_2$ & $2$ & $-$\\
\midrule
$1$ & $L_{-2},L_{-1}$ & $E$ & --- & ---\\
$1$ & $L_0$ & $V_1$ & $1$ & $+$\\
$1$ & $L_1$ & $U_1$ & $1$ & $+$\\
$1$ & $L_2$ & $U_0$ & $3$ & $-$\\
$1$ & $V_{-1}$ & $V_2$ & $1$ & $+$\\
$1$ & $U_0$ & $L_2$ & $13$ & $-$\\
$1$ & $V_0$ & $U_2$ & $1$ & $+$\\
$1$ & $V_1$ & $L_0$ & $2$ & $-$\\
$1$ & $V_2$ & $E$ & --- & ---\\
\midrule
$2$ & $L_{-2}$ & $E$ & --- & ---\\
$2$ & $L_{-1}$ & $V_1$ & $1$ & $+$\\
$2$ & $L_0$ & $U_1$ & $1$ & $+$\\
$2$ & $U_0$ & $L_1$ & $5$ & $-$\\
$2$ & $V_0$ & $V_2$ & $1$ & $+$\\
$2$ & $U_1$ & $L_0$ & $3$ & $+$\\
$2$ & $V_1$ & $U_2$ & $1$ & $+$\\
$2$ & $U_2$ & $L_{-1}$ & $1$ & $-$\\
$2$ & $V_2$ & $E$ & --- & ---\\
\end{longtable}
}

\begin{proof}
The complement calculation \eqref{eq:complement} and the outer
threshold prove that the start list is exhaustive. Each row then
follows from the recurrence \eqref{eq:recurrence}, using
\eqref{eq:values} and oddness at a boundary coordinate. Here $R$
alone is free. Reduce each affine coordinate modulo $6$ on its
fixed radius progression and continue until the first return or
escape. The intervening core inequalities are linear in $R$;
checking their direction and their value at the least allowed
radius proves them throughout the progression.

For the longest endpoint return, namely $s=1$ and
$U_0=(R,R+1)$, the second-coordinate list from time $0$ through
time $13$ is
\begin{align*}
 &R+1,\ -R+2,\ -R,\ R-3,\ R-1,\ -R+3,\ -R+2,\\
 &R-2,\ R-3,\ -R+1,\ -R+3,\ R,\ R-2,\ -R-1.
\end{align*}
The final pair is $-L_2$, and all intermediate second coordinates
lie in $[-R,R]$ already at $R=4$. This illustrates both the
negative sign and the first-return check.

For precision about the $E$ entries, when $s=0$ the paths are
\[
 V_0\ \xto{2,-}\ [-R,2]\longrightarrow E,\qquad
 U_2\ \xto{1,-}\ [-R-1,2]\longrightarrow E.
\]
Their displayed outer targets satisfy \eqref{eq:outer-threshold}.
When $s=1$, the path is
$V_2\xto{2,-}[-R+1,2]\longrightarrow E$.
When $s=2$, $V_2$ leaves $B$ in time $2$.
Every $L$ state marked $E$ leaves $B$ in time $1$.
Thus the compression of these transient paths to $E$ does not
discard any candidate cycle.
\end{proof}

\begin{lemma}[No endpoint aliases at small radii]\label{lem:noalias}
For every allowed $R$, the exceptional inner labels in
Table~\ref{tab:endpoints} are pairwise distinct and lie outside all
generic strips. The outer labels are distinct from the inner labels.
\end{lemma}
\begin{proof}
For $s=0$, the largest left inner coordinate is $-R+2$ and the
smallest right one is $R+1$, with $R\ge3$. For $s=1$, they are
$-R+2$ and $R$, with $R\ge4$. For $s=2$, they are $-R$ and
$R$, with $R\ge2$. Each gap is strictly positive already at the
least radius and increases with $R$. Different offsets on the
same side differ by a nonzero constant. Outer labels have second
coordinate $R+2$ rather than $R+1$. Finally,
\eqref{eq:complement} separates every inner endpoint from every
generic strip of its residue, and different residues cannot coincide.
\end{proof}
